\section{Alternative and Complementary Approaches}

\textcite[1652--1653]{jagannathan2003risk} distinguish between imposing structure on the covariance estimate, combining the sample matrix with a structured target and increasing the number of observations. This section considers factor models and principal components as structured estimators before examining portfolio constraints. The first two change the covariance estimator, whereas constraints act at the optimisation stage.

Factor models explain asset co-movement through a limited number of common factors and firm-specific shocks rather than estimating every pairwise relationship separately. Factors such as commodity prices, exchange rates or investor sentiment may cause shares with similar exposures to move together. For example, \textcite[186--187]{fan2007highdimensionalcovariancematrix} show that a three-factor model requires approximately \(4N\) parameters rather than the \(N(N+1)/2\) distinct parameters required by an unrestricted covariance matrix. The relevant factors must still be selected, and omitting an important factor may leave some co-movement unexplained. Factor estimators can therefore reduce sampling variance but introduce bias when their restrictions are inaccurate. Shrinkage limits this model risk by combining the structured estimator with the sample covariance matrix.

Principal-components analysis also reduces dimensionality but identifies common factors directly from patterns in the return data rather than selecting them beforehand. The first component explains the greatest amount of common variation, while subsequent components explain the greatest remaining variation. These components need not have clear economic interpretations. Using fewer components can reduce sampling noise, although discarded components may contain relevant covariance information. The number retained therefore remains a model-selection problem \parencite[191, 197--200]{bai2002determining}. By contrast, nonlinear shrinkage preserves all sample eigenvectors and adjusts their eigenvalues continuously.

Portfolio constraints reduce the effect of estimation error by restricting the weights that the optimiser can select rather than changing the covariance estimator. \textcite[1674--1677]{jagannathan2003risk} find that, after imposing non-negativity constraints, global minimum-variance portfolios using the sample covariance matrix perform almost as well as those using factor and shrinkage estimators. Their simulations show that this result continues to hold when returns are strongly influenced by a common factor. Constraints may therefore reduce the performance advantage of shrinkage, although they can also exclude useful positions and reduce the performance of estimators that are already stable.

The role of portfolio constraints also helps explain the empirical design below. By using unconstrained weights, the JSE application allows the effects of the different covariance estimators on short positions and portfolio concentration to remain visible.
